%
% A proof-by-cases environment. This gives you a nicely-indented list
% of cases for use in proofs.
%
% Example:
%
% Case 1 (x >= 0): herp.
%
% Case 2 (x < 0): derp.
%
\ifx\havemjoproofbycases\undefined
\def\havemjoproofbycases{1}


% Used below to define pcases. The ``loadonly'' parameter prevents
% a very bad interaction with the beamer document class.
\usepackage[loadonly]{enumitem}

% Needed to perform division in the definition of \singleblskip.
\usepackage{calc}

% A \baselineskip without the \baselinestretch scaling factor. Even
% though \baselinestretch defaults to 1.0, it doesn't really have that
% value. Thus to avoid division by zero, we need to do the ``is this
% thing empty?'' hack.
%
% If we use \baselineskip instead of \singleblskip in our list, things
% get real ugly when the text is e.g. double-spaced.
%
\newlength{\singleblskip}
\setlength{\singleblskip}{
  \baselineskip/\real{
      \if\relax\detokenize{\baselinestretch}\relax
        \baselinestretch%
      \else
        1%
      \fi}
}

% Using the enumitem package, we define a new type of list, called
% ``pcases'' (proof by cases). Each case has a label with an arabic
% numeral (the case number), but also a \thiscase identifier. The
% macro \thiscase is defined below by the \case command, and gives the
% name or conditions or whatever that distinguish one case from
% another.
\newlist{pcases}{enumerate}{1}
\setlist[pcases]{
  label=\textit{Case~\arabic*}:~\protect\thiscase\textit{.},
  ref=\arabic*,
  align=left,
  listparindent=\parindent,
  parsep=\parskip,
  itemsep=\singleblskip}

% The optional argument here gets stuffed into the \thiscase macro, to
% be called by pcases when it creates this list item. The \hfill is a
% hack intended to force the proof to start on a new line, rather than
% right after the colon. A \newline where the \hfill is does not work,
% so we consume the rest of the line instead.
\newcommand{\case}[1][]{
  \def\thiscase{#1}%
  \item \hfill\par\vspace{\singleblskip}
}


\fi
